%% chapter8.tex – Future Scope
\chapter{Future Scope}
\label{chap:future}

The WaveGuard~V4 prototype has demonstrated core feasibility.
The following enhancements are planned for subsequent iterations.

\section{Multi-Buoy Mesh Network}

Deploying a network of 5--10 buoys at 1~km spacing along the Kanayannur coastline
would enable directional wave propagation tracking.
ESP32 devices support ESP-NOW (a proprietary peer-to-peer protocol at 1~Mbps),
allowing inter-buoy communication without a WiFi router.
A gateway buoy at shore would aggregate all readings before forwarding to the backend.

\section{Machine-Learning Based Anomaly Detection}

The current threshold-based classifier can produce false positives during
boat-induced wake events that share an RMS footprint similar to swell.
An LSTM autoencoder trained on labelled wave spectra could distinguish
\textit{Kallakkadal}-type long-period swell from boat wake using frequency-domain
features derived from the 50~Hz time series.
The ESP32's 520~KB RAM is sufficient for inference with TensorFlow Lite for
Microcontrollers (TFLite-Micro) for models $<$50~KB.

\section{Solar Power Enhancement}

Replacing the single 18650 cell with a 6~V / 3~W solar panel and a 3-cell 18650
pack managed by a CN3791 MPPT charger would allow indefinite deployment.
A night-time deep-sleep schedule (ESP32 ULP coprocessor at 10~µA) would further
reduce average consumption to $<$30~mA, extending useful solar coverage to cloudy days.

\section{Mobile Application}

A React Native mobile app would extend alert delivery beyond SMS to push
notifications (FCM/APNs), enabling real-time notifications even when cellular SMS
credits are exhausted.
The existing FastAPI REST backend requires only a JWT-secured push-token endpoint to
support mobile clients.

\section{Enhanced Notification Channels}

\begin{itemize}
  \item \textbf{SMTP email alerts} via GMail OAuth2 API for fisheries department
        officials.
  \item \textbf{Telegram Bot API} for community group broadcasts.
  \item \textbf{WhatsApp Business API} (Twilio/360dialog) for broader reach.
\end{itemize}

\section{Real-Time Cloud Integration}

Migrating the SQLite backend to a cloud-hosted TimescaleDB (PostgreSQL time-series
extension) would allow:
\begin{itemize}
  \item Multi-buoy data aggregation from a central dashboard.
  \item Long-term historical analysis for \textit{Kallakkadal} season profiling.
  \item Public API for research institutions.
\end{itemize}
ESP32 cellular data upload via SIM800L TCP/IP stack can replace the shore-based WiFi
relay, removing the 50~m range constraint.

\section{Advanced Data Visualisation}

A Grafana dashboard connected to TimescaleDB would provide:
\begin{itemize}
  \item Heat-map of alert frequency along the 1~km transect.
  \item Spectral power density plot from FFT of the 50~Hz time series.
  \item Correlation overlays with Open-Meteo wave height forecast data.
\end{itemize}

\section{Waterproofing and Buoy Mechanical Redesign}

The current ABS enclosure is tested to IP54; upgrading to a moulded HDPE sphere
(IP68, 10~m depth rating) with a through-hull cable gland for the GPS antenna would
support the 5--10~m deployment depth specified in the project brief.
A gyroscope-stabilised mounting platform inside the sphere would reduce the
false-positive rate from rotational buoy pitch/roll.
